\section{Intelligence as a Second Attractor}

This section treats intelligence not as a direct extension of life, but as a second stable structure supported by the life attractor. This view separates the conditions for life from the conditions for intelligence.

\subsection{Position of Intelligence}

Intelligence is often treated as an advanced function of living systems. However, life is already an attractor structure. Many living systems persist for long periods within that attractor without developing intelligence. Intelligence is therefore not merely a strengthening of the life attractor.

\subsection{Second Attractor}

Intelligence appears when information is not only flowing, but also retained, referenced, and reused. Past states influence present states, and present states modify future trajectories. This creates recurrent loops of information flow.

Such loops stabilize under conditions that differ from those required for life alone. Intelligence is therefore interpreted as a second attractor supported by, but not identical to, the life attractor.

\subsection{Structural Independence}

Formally, one may write
\begin{equation}
  \Omega_{\mathrm{intel}} \subset \Omega_{\mathrm{life}},
\end{equation}
but the internal structure of \(\Omega_{\mathrm{intel}}\) is determined by additional constraints. These include persistent reference to state history, temporal accumulation of information, and sharing of information across multiple components. Such conditions are not required for biological maintenance alone.

\subsection{Probability Structure}

Following Section 8, the emergence of intelligence is treated as the satisfaction of conditions distinct from the emergence of life. Its probability therefore has a conditional structure. A schematic factorization is
\begin{equation}
  P_{\mathrm{intel}} \approx P_{\mathrm{life}}\,P_{\mathrm{second}\mid\mathrm{life}},
\end{equation}
where \(P_{\mathrm{second}\mid\mathrm{life}}\) denotes the probability of reaching the second attractor given the life attractor. This expression is not a numerical estimate. It records the two-stage structure of the problem.

\subsection{Non-Teleological Interpretation}

Treating intelligence as a second attractor does not give it a privileged status. An attractor is not a goal. It is a structure that persists when its conditions are satisfied. Intelligence is therefore neither inevitable nor a final form. It is one possible stable structure.

\subsection{Relation to Internal Time}

In the notation of Section 7, life is a structure in which \(\tau\) continues to progress. Intelligence is characterized by the reference and reuse of the history of \(\tau\). In this sense, intelligence is not merely flow. It is recursive information flow.

\subsection{Summary}

Intelligence is not treated as a simple extension of life. It is a stable structure of recursive information flow supported by the life attractor. This interpretation explains why intelligence has additional conditions and is not universal among living systems.

\subsection{Connection to Further Work}

The next question is how this second attractor is implemented. Its realization may involve memory, temporal integration, language, and inter-individual sharing. These issues are beyond the minimal biological framework developed here.
